%% toc: P1 8 8.1,8.2,8.5,8.6
\section{Proofs for notation, binding and functional completeness}
\label{sec:annex-p01-w02-p01-024-notation}

\begin{annexitem}{annex:p01-w02-p01-024-fv-occurrence}{thm:p01-w02-p01-024-fv-occurrence}{The free-variable set collects the free occurrences}
\begin{proof}
Write \(\operatorname{OV}(\varphi)=\{\varphi[a]:a\in\operatorname{VO}(\varphi)\}\)
and \(\operatorname{FO}(\varphi)=\{\varphi[a]:a\in\operatorname{VO}(\varphi),
\ a\text{ free in }\varphi\}\).

\emph{Terms.}  For every term \(t\),
\(\{t[a]:a\in\operatorname{VO}(t)\}=\operatorname{Var}(t)\).  For a
variable both sides are \(\{x\}\).  For \(f(t_1,\ldots,t_n)\) the left side
is \(\bigcup_{i=1}^n\{t_i[a]:a\in\operatorname{VO}(t_i)\}\), which by the
induction hypothesis is \(\bigcup_{i=1}^n\operatorname{Var}(t_i)\); for
\(n=0\) both sides are empty.

\emph{A criterion for being bound.}  For \(a\in\operatorname{VO}(\varphi)\)
let \(P_\varphi(a)\) be the set of \(b\in\operatorname{QO}(\varphi)\) with
\(b\sqsubseteq a\) that quantify \(\varphi[a]\).  If some \(b\) binds \(a\),
then \(b\in P_\varphi(a)\).  Conversely, the members of \(P_\varphi(a)\) are
prefixes of the one address \(a\), hence finitely many and linearly ordered
by \(\sqsubseteq\); if \(P_\varphi(a)\neq\varnothing\), its longest member
\(b\) binds \(a\), since no \(b'\neq b\) with \(b\sqsubseteq b'\sqsubseteq a\)
lies in \(P_\varphi(a)\).  So \(a\) is bound iff
\(P_\varphi(a)\neq\varnothing\).

\emph{Compositional clauses.}  Let \(\varphi\) be compound with immediate
subformula \(\psi\) at index \(i\) (so \(i=1\) for \(\neg\psi\) and for
\(Qx\,\psi\), and \(i\in\{1,2\}\) for \(\psi_1\circ\psi_2\)), and let
\(a=i\cdot a'\) with \(a'\in\operatorname{VO}(\psi)\).  By the definitions,
\(\varphi[a]=\psi[a']\), and the members of \(\operatorname{QO}(\varphi)\)
other than \(\epsilon\) that are prefixes of \(a\) are exactly the addresses
\(i\cdot b\) with \(b\in\operatorname{QO}(\psi)\) and \(b\sqsubseteq a'\);
moreover \(\varphi\{i\cdot b\}=\psi\{b\}\), so \(i\cdot b\) quantifies the
same variable in \(\varphi\) as \(b\) does in \(\psi\).
\begin{enumerate}
\item If \(\varphi\) is \(\neg\psi\) or \(\psi_1\circ\psi_2\), then
  \(\epsilon\notin\operatorname{QO}(\varphi)\), so
  \(P_\varphi(i\cdot a')=\{i\cdot b:b\in P_\psi(a')\}\): the occurrence
  \(i\cdot a'\) is free in \(\varphi\) iff \(a'\) is free in \(\psi\).
\item If \(\varphi=Qx\,\psi\), then \(\epsilon\in\operatorname{QO}(\varphi)\)
  quantifies \(x\) and is a prefix of every address.  The occurrence \(0\)
  has \(\varphi[0]=x\), so \(\epsilon\in P_\varphi(0)\) and \(0\) is bound.
  For \(a=1\cdot a'\): if \(\psi[a']=x\), then \(\epsilon\in P_\varphi(a)\)
  and \(a\) is bound; if \(\psi[a']\neq x\), then \(\epsilon\notin
  P_\varphi(a)\) and \(P_\varphi(a)=\{1\cdot b:b\in P_\psi(a')\}\), so
  \(a\) is free iff \(a'\) is free in \(\psi\).
\end{enumerate}

\emph{Induction on \(\varphi\).}  For \(r=t\) and \(R(t_1,\ldots,t_n)\),
\(\operatorname{QO}=\varnothing\), so every occurrence is free and
\(\operatorname{OV}=\operatorname{FO}=\bigcup_i\operatorname{Var}(t_i)\) by
the clause for terms; this is both \(\operatorname{Var}\) and \(\FV\) of the
atomic formula.  For \(\bottom\) all four sets are empty.  For \(\neg\psi\)
and \(\psi_1\circ\psi_2\), clause 1 gives
\(\operatorname{FO}(\neg\psi)=\operatorname{FO}(\psi)\) and
\(\operatorname{FO}(\psi_1\circ\psi_2)=\operatorname{FO}(\psi_1)\cup
\operatorname{FO}(\psi_2)\); the same unions hold for
\(\operatorname{OV}\) directly from the definition of
\(\operatorname{VO}\); and \(\FV\) and \(\operatorname{Var}\) satisfy the
same recursion.  For \(Qx\,\psi\), clause 2 gives
\(\operatorname{FO}(Qx\,\psi)=\{\psi[a']:a'\text{ free in }\psi,\
\psi[a']\neq x\}=\operatorname{FO}(\psi)\setminus\{x\}\), which by the
induction hypothesis is \(\FV(\psi)\setminus\{x\}=\FV(Qx\,\psi)\); and
\(\operatorname{OV}(Qx\,\psi)=\{x\}\cup\operatorname{OV}(\psi)=\{x\}\cup
\operatorname{Var}(\psi)\), which is \(\operatorname{Var}(Qx\,\psi)\)
because the quantifier clause for \(\operatorname{Var}\) adds the bound
variable.

\emph{Consequences.}  By the displayed identity, \(x\in\FV(\varphi)\) iff
some occurrence of \(x\) is free, so \(x\notin\FV(\varphi)\) iff every
occurrence of \(x\) is bound.  A sentence is a formula with
\(\FV(\varphi)=\varnothing\), which holds iff no variable has a free
occurrence, that is, iff every variable occurrence is bound.
\end{proof}
\end{annexitem}

\begin{annexitem}{annex:p01-w02-p01-024-positive-diagram}{thm:p01-w02-p01-024-positive-diagram}{What each diagram characterizes}
\begin{proof}
Write \(\mathcal M_M\) for the \(\Lang(M)\)-expansion of
\(\mathcal M\) with \(c_a^{\mathcal M_M}=a\).  An \(\Lang(M)\)-expansion
\(\mathcal N'\) of \(\mathcal N\) interprets every symbol of \(\Lang\) as
\(\mathcal N\) does and assigns an element \(c_a^{\mathcal N'}\) to each
\(c_a\).

\emph{Closed terms.}  Let \(h:M\to N\) be a function that preserves the
constants and functions of \(\Lang\) (clauses 1 and 2 of the definition of
homomorphism), and let \(\mathcal N'\) interpret \(c_a\) as \(h(a)\).  Then
\(t^{\mathcal N'}=h(t^{\mathcal M_M})\) for every closed
\(\Lang(M)\)-term \(t\).  For \(t=c_a\) both sides are \(h(a)\); for a
constant \(c\) of \(\Lang\) the claim is \(c^{\mathcal N}=h(c^{\mathcal M})\);
for \(t=f(t_1,\ldots,t_n)\),
\(h(f^{\mathcal M}(t_1^{\mathcal M_M},\ldots,t_n^{\mathcal M_M}))
=f^{\mathcal N}(h(t_1^{\mathcal M_M}),\ldots,h(t_n^{\mathcal M_M}))
=f^{\mathcal N}(t_1^{\mathcal N'},\ldots,t_n^{\mathcal N'})\) by
preservation of \(f\) and the induction hypothesis.

\emph{Part 1, forward.}  Let \(h\) be a homomorphism and \(\mathcal N'\) as
above.  Let \(\theta\in\operatorname{Diag}^+(\mathcal M)\).  If \(\theta\)
is \(s=t\), then \(s^{\mathcal M_M}=t^{\mathcal M_M}\), so
\(s^{\mathcal N'}=h(s^{\mathcal M_M})=h(t^{\mathcal M_M})=t^{\mathcal N'}\).
If \(\theta\) is \(R(t_1,\ldots,t_n)\), then
\((t_1^{\mathcal M_M},\ldots,t_n^{\mathcal M_M})\in R^{\mathcal M}\), so
by clause 3 of the definition of homomorphism
\((t_1^{\mathcal N'},\ldots,t_n^{\mathcal N'})\in R^{\mathcal N}\).  Hence
\(\mathcal N'\entails\theta\).

\emph{Part 1, backward.}  Let
\(\mathcal N'\entails\operatorname{Diag}^+(\mathcal M)\) and put
\(h(a)=c_a^{\mathcal N'}\).  For a constant \(c\) of \(\Lang\), the
sentence \(c=c_{c^{\mathcal M}}\) is atomic and true in \(\mathcal M_M\), so
\(c^{\mathcal N}=c_{c^{\mathcal M}}^{\mathcal N'}=h(c^{\mathcal M})\).  For
an \(n\)-ary \(f\) and \(a_1,\ldots,a_n\in M\) with
\(b=f^{\mathcal M}(a_1,\ldots,a_n)\), the sentence
\(f(c_{a_1},\ldots,c_{a_n})=c_b\) is atomic and true in \(\mathcal M_M\), so
\(f^{\mathcal N}(h(a_1),\ldots,h(a_n))=h(b)\).  If
\((a_1,\ldots,a_n)\in R^{\mathcal M}\), then
\(R(c_{a_1},\ldots,c_{a_n})\in\operatorname{Diag}^+(\mathcal M)\), so
\((h(a_1),\ldots,h(a_n))\in R^{\mathcal N}\).  Thus \(h\) is a
homomorphism.

\emph{Part 2.}  Since
\(\operatorname{Diag}(\mathcal M)\supseteq\operatorname{Diag}^+(\mathcal M)\),
the backward argument of part 1 applies to an expansion satisfying
\(\operatorname{Diag}(\mathcal M)\) and gives a homomorphism
\(h(a)=c_a^{\mathcal N'}\).  For distinct \(a,b\in M\) the sentence
\(\neg(c_a=c_b)\) is the negation of an atomic sentence false in
\(\mathcal M_M\), so it lies in \(\operatorname{Diag}(\mathcal M)\) and
\(h(a)\neq h(b)\).  If \((a_1,\ldots,a_n)\notin R^{\mathcal M}\), then
\(\neg R(c_{a_1},\ldots,c_{a_n})\in\operatorname{Diag}(\mathcal M)\), so
\((h(a_1),\ldots,h(a_n))\notin R^{\mathcal N}\); with preservation this is
reflection.  So \(h\) is an embedding.  Conversely, if \(h\) is an embedding
and \(\mathcal N'\) interprets \(c_a\) as \(h(a)\), part 1 gives
\(\mathcal N'\entails\operatorname{Diag}^+(\mathcal M)\).  A negated atomic
member \(\neg(s=t)\) has \(s^{\mathcal M_M}\neq t^{\mathcal M_M}\), and
injectivity with the closed-term identity gives
\(s^{\mathcal N'}\neq t^{\mathcal N'}\); a negated atomic member
\(\neg R(t_1,\ldots,t_n)\) has
\((t_1^{\mathcal M_M},\ldots,t_n^{\mathcal M_M})\notin R^{\mathcal M}\),
and reflection gives
\((t_1^{\mathcal N'},\ldots,t_n^{\mathcal N'})\notin R^{\mathcal N}\).
Hence \(\mathcal N'\entails\operatorname{Diag}(\mathcal M)\).
\end{proof}
\end{annexitem}

\begin{annexitem}{annex:p01-w02-p01-024-nand-nor}{thm:p01-w02-p01-024-nand-nor}{NAND and NOR are complete for positive arities}
\begin{proof}
Formulas are evaluated by the recursive truth clauses, and the truth
function of a formula in the variables \(p_1,\ldots,p_n\) is well defined
by coincidence (\cref{sec:prop-syntax-semantics-04}).

\emph{Step 1: \(\{\neg,\land\}\) is positive-arity complete.}  Let
\(n\geq1\) and \(f:\{0,1\}^n\to\{0,1\}\).  If \(f\) is constantly \(0\),
the formula \(p_1\land\neg p_1\) has value \(0\) under every valuation and
its only variable is among \(p_1,\ldots,p_n\).  Otherwise let
\(A=\{\bar a:f(\bar a)=1\}\neq\varnothing\).  For \(\bar a\in A\) let
\(\lambda_{\bar a}\) be the left-associated conjunction
\(\ell_1\land\cdots\land\ell_n\), where \(\ell_i=p_i\) if \(a_i=1\) and
\(\ell_i=\neg p_i\) if \(a_i=0\).  A conjunction has value \(1\) iff each
conjunct does, so \(\lambda_{\bar a}\) has value \(1\) under a valuation
with \(v(p_i)=b_i\) iff \(\bar b=\bar a\).  Define
\(\alpha\mathbin{\dot\lor}\beta\defeq\neg(\neg\alpha\land\neg\beta)\); its
value is \(1\) iff the value of \(\alpha\) or of \(\beta\) is \(1\).  Let
\(\delta\) be the \(\dot\lor\)-combination, in any fixed order, of the
formulas \(\lambda_{\bar a}\) for \(\bar a\in A\).  Under a valuation with
\(v(p_i)=b_i\), \(\delta\) has value \(1\) iff some \(\lambda_{\bar a}\)
does, iff \(\bar b\in A\), iff \(f(\bar b)=1\).  So \(\delta\) has truth
function \(f\).

\emph{Step 2: translation into NAND.}  Define \(\tau(p)=p\),
\(\tau(\neg\alpha)=\operatorname{NAND}(\tau\alpha,\tau\alpha)\) and
\(\tau(\alpha\land\beta)=\operatorname{NAND}(\operatorname{NAND}(\tau\alpha,
\tau\beta),\operatorname{NAND}(\tau\alpha,\tau\beta))\).  Since
\(\operatorname{NAND}(x,x)=1-x\) and
\(\operatorname{NAND}(\operatorname{NAND}(x,y),\operatorname{NAND}(x,y))
=1-\operatorname{NAND}(x,y)=xy\), induction on \(\alpha\) shows that
\(\tau\alpha\) has the same value as \(\alpha\) under every valuation and
the same variables.  Applied to the formula of step 1, this represents
\(f\) by a formula over \(\{\operatorname{NAND}\}\).

\emph{Step 3: translation into NOR.}  Define \(\tau'(p)=p\),
\(\tau'(\neg\alpha)=\operatorname{NOR}(\tau'\alpha,\tau'\alpha)\) and
\(\tau'(\alpha\land\beta)=\operatorname{NOR}(\operatorname{NOR}(\tau'\alpha,
\tau'\alpha),\operatorname{NOR}(\tau'\beta,\tau'\beta))\).  Since
\(\operatorname{NOR}(x,x)=1-x\) and \(\operatorname{NOR}(1-x,1-y)=1\) iff
\(1-x=0\) and \(1-y=0\), iff \(x=y=1\), the same induction applies.

\emph{Step 4: the constant \(\bottom\).}  Over
\(\{\operatorname{NAND},\bottom\}\) the closed formulas \(\bottom\) and
\(\operatorname{NAND}(\bottom,\bottom)\) have values \(0\) and \(1\), so both
nullary functions are represented; positive arities are covered by step 2.
Over \(\{\operatorname{NOR},\bottom\}\) the closed formulas are \(\bottom\)
and \(\operatorname{NOR}(\bottom,\bottom)\), and step 3 covers positive
arities.
\end{proof}
\end{annexitem}

\begin{annexitem}{annex:p01-w02-p01-024-no-closed}{thm:p01-w02-p01-024-no-closed}{No nullary representation without a constant}
\begin{proof}
\emph{Part 1.}  Every formula over \(C\) is a propositional variable or has
the form \(c(\varphi_1,\ldots,\varphi_k)\) for a connective \(c\in C\) of
arity \(k\).  As \(C\) has no nullary member, \(k\geq1\), and by the
induction hypothesis \(\varphi_1\) contains a variable, so the compound
formula does too.  A representation of a nullary function is a formula
whose variables are among the empty list \(p_1,\ldots,p_0\), that is, a
closed formula; there is none.  So neither nullary function is
represented, and \(C\) is not nullary-inclusive complete.  NAND and NOR are
binary, so this applies to \(\{\operatorname{NAND}\}\) and
\(\{\operatorname{NOR}\}\).

\emph{Part 2.}  If \(C\) is nullary-inclusive complete, it is positive-arity
complete by definition, and it has a closed formula, so by part 1 it has a
nullary member.  Conversely, let \(C\) be positive-arity complete with a
nullary member \(k\) whose value is \(b\).  Positive-arity completeness for
\(n=1\) gives a formula \(\nu\) in the variable \(p_1\) whose truth
function is \(x\mapsto1-x\).  Let \(\nu[k/p_1]\) be the result of replacing
every occurrence of \(p_1\) in \(\nu\) by \(k\).  By induction on \(\nu\),
the value of \(\nu[k/p_1]\) under a valuation \(v\) equals the value of
\(\nu\) under the valuation that agrees with \(v\) except that it sends
\(p_1\) to \(b\); since \(\nu\) contains no other variable, that value is
\(1-b\).  So the closed formulas \(k\) and \(\nu[k/p_1]\) represent the two
nullary functions, and positive arities are covered by hypothesis.
\end{proof}
\end{annexitem}

\begin{annexitem}{annex:p01-w02-p01-024-post}{thm:p01-w02-p01-024-post}{Post's criterion for positive arities}
\begin{proof}
All functions below have positive arity.  Write \(\neg\) for the unary
function \(x\mapsto1-x\), \(\mathrm{id}=e^1_1\), \(k_0,k_1\) for the unary
constant functions with values \(0,1\), \(\land\) and \(\to\) for the binary
functions of rows 0001 and 1101 of Table~\ref{tab:p01-sixteen}, \(\oplus\)
for the binary function of row 0110, and \(\operatorname{maj}\) for the
ternary majority function.  For \(f\in O_n\) and \(1\leq i\leq n\),
\(\bar a\oplus e_i\) is \(\bar a\) with coordinate \(i\) changed.

\emph{Clone facts.}  (i) If \(K\) is a clone, \(f\in K\cap O_m\) and
\(g_1,\ldots,g_m\in K\cap O_n\), then \(f(g_1,\ldots,g_m)\in K\); in
particular identifying all arguments, \(x\mapsto f(x,\ldots,x)\), is
\(f(\mathrm{id},\ldots,\mathrm{id})\in K\), and a unary \(u\in K\) yields
the binary functions \(u(e^2_1)\) and \(u(e^2_2)\) in \(K\).  (ii) If
\(C\subseteq X\) for a clone \(X\), then \(\langle C\rangle\subseteq X\),
because \(X\) is one of the clones intersected.

\emph{Lemma A (affine counting).}  A nonconstant affine
\(f=c_0\oplus c_1a_1\oplus\cdots\oplus c_na_n\) takes the value \(1\) at
exactly \(2^{n-1}\) inputs.  Pick \(k\) with \(c_k=1\); the map
\(\bar a\mapsto\bar a\oplus e_k\) is a bijection of \(\{0,1\}^n\) onto
itself that changes the value of \(f\), so it maps the inputs with value
\(1\) bijectively onto those with value \(0\).  For \(n=2\) an affine
function therefore takes the value \(1\) at \(0\), \(2\) or \(4\) inputs.

\emph{Lemma B (affinity test).}  \(f\in O_n\) is affine iff
\(f(\bar b)\oplus f(\bar b\oplus e_i)\oplus f(\bar b\oplus e_j)\oplus
f(\bar b\oplus e_i\oplus e_j)=0\) for all \(\bar b\) and all \(i\neq j\).
If \(f\) is affine, each \(c_i\) term occurs an even number of times in the
sum and \(c_0\) four times, so the sum is \(0\).  Conversely, put
\(c_0=f(\bar0)\) and \(c_i=f(e_i)\oplus f(\bar0)\), and prove
\(f(\bar a)=c_0\oplus\bigoplus_ic_ia_i\) by induction on the number of
coordinates of \(\bar a\) equal to \(1\).  For zero or one such coordinate
this is the definition.  Otherwise pick \(i\neq j\) with \(a_i=a_j=1\) and
put \(\bar b=\bar a\oplus e_i\oplus e_j\).  The hypothesis at \(\bar b\)
gives \(f(\bar a)=f(\bar b)\oplus f(\bar b\oplus e_i)\oplus
f(\bar b\oplus e_j)\), and each argument on the right has fewer coordinates
equal to \(1\).  By the induction hypothesis the right side is
\(3c_0\oplus3\bigoplus_kc_kb_k\oplus c_i\oplus c_j=c_0\oplus
\bigoplus_kc_ka_k\), since \(\bar a\) differs from \(\bar b\) exactly at
\(i\) and \(j\), where \(b_i=b_j=0\).

\emph{Part 1: the five classes are clones.}  Projections:
\(e^n_i(\bar0)=0\), \(e^n_i(\bar1)=1\),
\(e^n_i(\bar a^{\,c})=1-a_i\), \(e^n_i\) is monotone, and
\(e^n_i=a_i\) is affine.  Composition: let \(h=f(g_1,\ldots,g_m)\).  If
\(f,g_j\in\mathsf T_0\), then \(h(\bar0)=f(0,\ldots,0)=0\); likewise for
\(\mathsf T_1\).  If \(f,g_j\in\mathsf S\), then
\(h(\bar a^{\,c})=f(1-g_1(\bar a),\ldots,1-g_m(\bar a))
=1-h(\bar a)\).  If \(f,g_j\in\mathsf M\) and \(\bar a\leq\bar b\), then
\(g_j(\bar a)\leq g_j(\bar b)\) for every \(j\), so
\(h(\bar a)\leq h(\bar b)\).  If \(f=c_0\oplus\bigoplus_jc_jx_j\) and
\(g_j=d_{j0}\oplus\bigoplus_id_{ji}a_i\), then, computing in the
two-element field,
\(h(\bar a)=(c_0\oplus\bigoplus_jc_jd_{j0})\oplus\bigoplus_i
(\bigoplus_jc_jd_{ji})a_i\), which is affine.

\emph{Part 1: memberships of seven witnesses.}
\(k_0\in\mathsf T_0\cap\mathsf M\cap\mathsf L\) and
\(k_0\notin\mathsf T_1\cup\mathsf S\) (\(k_0(1)=0\), and
\(k_0(1-a)=0\neq1-k_0(a)\)).  Symmetrically
\(k_1\in\mathsf T_1\cap\mathsf M\cap\mathsf L\) and
\(k_1\notin\mathsf T_0\cup\mathsf S\).  \(\neg\in\mathsf S\cap\mathsf L\)
(\(\neg(1-a)=a\), and \(\neg a=1\oplus a\)) and
\(\neg\notin\mathsf T_0\cup\mathsf T_1\cup\mathsf M\) (\(\neg0=1\),
\(\neg1=0\), and \(0\leq1\) with \(\neg0>\neg1\)).
\(\land\in\mathsf T_0\cap\mathsf T_1\cap\mathsf M\); \(\land\notin\mathsf S\)
because \(\land(1,0)=0\neq1-\land(0,1)\); \(\land\notin\mathsf L\) by
Lemma A, as it takes the value \(1\) at one input.
\(\oplus\in\mathsf T_0\) and \(\oplus\notin\mathsf M\), since
\((0,1)\leq(1,1)\) while \(0\oplus1=1>0=1\oplus1\).
\(\to\in\mathsf T_1\) and \(\to\notin\mathsf M\), since
\((0,0)\leq(1,0)\) while \((0\to0)=1>0=(1\to0)\).
\(\operatorname{maj}\in\mathsf S\), because at least two of
\(1-a,1-b,1-c\) equal \(1\) iff at most one of \(a,b,c\) does; and
\(\operatorname{maj}\notin\mathsf L\): if
\(\operatorname{maj}=c_0\oplus c_1a\oplus c_2b\oplus c_3c\), the values at
\((0,0,0)\), \((1,0,0)\), \((0,1,0)\) and \((0,0,1)\), all \(0\), force
\(c_0=c_1=c_2=c_3=0\), contradicting
\(\operatorname{maj}(1,1,1)=1\).

\emph{Part 1: properness and non-inclusion.}  Each class omits a witness
(\(\neg\) is outside \(\mathsf T_0,\mathsf T_1,\mathsf M\); \(\land\) is
outside \(\mathsf S,\mathsf L\)), so each is proper.  For distinct classes
\(X,Y\) a member of \(X\setminus Y\) is:
\(k_0\) for \((\mathsf T_0,\mathsf T_1)\), \((\mathsf T_0,\mathsf S)\),
\((\mathsf M,\mathsf T_1)\), \((\mathsf M,\mathsf S)\),
\((\mathsf L,\mathsf T_1)\), \((\mathsf L,\mathsf S)\);
\(k_1\) for \((\mathsf T_1,\mathsf T_0)\), \((\mathsf T_1,\mathsf S)\),
\((\mathsf M,\mathsf T_0)\), \((\mathsf L,\mathsf T_0)\);
\(\oplus\) for \((\mathsf T_0,\mathsf M)\);
\(\land\) for \((\mathsf T_0,\mathsf L)\), \((\mathsf T_1,\mathsf L)\),
\((\mathsf M,\mathsf L)\);
\(\to\) for \((\mathsf T_1,\mathsf M)\);
\(\neg\) for \((\mathsf S,\mathsf T_0)\), \((\mathsf S,\mathsf T_1)\),
\((\mathsf S,\mathsf M)\), \((\mathsf L,\mathsf M)\);
\(\operatorname{maj}\) for \((\mathsf S,\mathsf L)\).
These are all twenty ordered pairs, so no class is included in another.

\emph{Part 2, necessity.}  If \(C\subseteq X\) for one of the five classes,
then \(\langle C\rangle\subseteq X\neq O\) by clone fact (ii) and part 1.

\emph{Part 2, sufficiency.}  Let \(f_0,f_1,f_s,f_m,f_l\in C\) lie outside
\(\mathsf T_0,\mathsf T_1,\mathsf S,\mathsf M,\mathsf L\) respectively, and
let \(K=\langle C\rangle\).

\emph{Step A.}  \(u_0(x)=f_0(x,\ldots,x)\in K\) has \(u_0(0)=1\), so
\(u_0=\neg\) or \(u_0=k_1\).  \(u_1(x)=f_1(x,\ldots,x)\in K\) has
\(u_1(1)=0\), so \(u_1=\neg\) or \(u_1=k_0\).  Hence \(\neg\in K\) or
\(k_0,k_1\in K\).

\emph{Step B: both constants.}  Suppose \(\neg\in K\).  As
\(f_s\notin\mathsf S\), some \(\bar a\) has
\(f_s(\bar a^{\,c})=f_s(\bar a)\).  Let \(h_i=\mathrm{id}\) if \(a_i=1\) and
\(h_i=\neg\) if \(a_i=0\), and let \(u=f_s(h_1,\ldots,h_n)\in K\).  Then
\(h_i(1)=a_i\) and \(h_i(0)=1-a_i\), so \(u(1)=f_s(\bar a)\) and
\(u(0)=f_s(\bar a^{\,c})=u(1)\): \(u\) is a constant \(k_b\), and
\(\neg(k_b)=k_{1-b}\in K\).  With Step A, \(k_0,k_1\in K\) in every case.

\emph{Step C: negation.}  As \(f_m\notin\mathsf M\), there are
\(\bar a\leq\bar b\) with \(f_m(\bar a)=1\) and \(f_m(\bar b)=0\).  Changing
the coordinates where \(\bar a\) and \(\bar b\) differ one at a time gives
a chain \(\bar a=\bar d^0\leq\bar d^1\leq\cdots\leq\bar d^r=\bar b\), each
step changing one coordinate from \(0\) to \(1\); since the values at the
ends are \(1\) and \(0\), some step has \(f_m(\bar d^{s-1})=1\) and
\(f_m(\bar d^s)=0\), and the two tuples differ only at one coordinate
\(k\).  Let \(h_k=\mathrm{id}\) and \(h_i=k_{d^{s-1}_i}\) for \(i\neq k\).
Then \(u=f_m(h_1,\ldots,h_n)\in K\) has \(u(0)=f_m(\bar d^{s-1})=1\) and
\(u(1)=f_m(\bar d^s)=0\), so \(\neg=u\in K\).

\emph{Step D: conjunction.}  By Lemma B there are \(\bar b\) and \(i\neq j\)
with \(f_l(\bar b)\oplus f_l(\bar b\oplus e_i)\oplus f_l(\bar b\oplus e_j)
\oplus f_l(\bar b\oplus e_i\oplus e_j)=1\).  Define binary
\(g(x,y)=f_l(h_1,\ldots,h_n)\), where \(h_i\) is \(e^2_1\) if \(b_i=0\) and
\(\neg(e^2_1)\) if \(b_i=1\), \(h_j\) is \(e^2_2\) if \(b_j=0\) and
\(\neg(e^2_2)\) if \(b_j=1\), and \(h_k=k_{b_k}(e^2_1)\) otherwise; all
lie in \(K\) by clone fact (i).  Then
\(g(x,y)=f_l(\bar b\oplus xe_i\oplus ye_j)\), so \(g\in K\) takes the value
\(1\) at an odd number of its four inputs.  If at three, replace \(g\) by
\(\neg(g)\in K\), which takes it at exactly one input \((u,w)\).  Then
\(g'(x,y)=g(h(x),h'(y))\), with \(h=\mathrm{id}\) if \(u=1\) and \(\neg\) if
\(u=0\) and \(h'\) chosen likewise from \(w\), takes the value \(1\) exactly
at \((1,1)\); so \(g'=\land\in K\).

\emph{Step E: every function.}  For a formula \(\varphi\) over
\(\{\neg,\land\}\) in the variables \(p_1,\ldots,p_n\), its truth function
\(f_\varphi\in O_n\) lies in \(K\): \(f_{p_i}=e^n_i\),
\(f_{\neg\alpha}=\neg(f_\alpha)\) and
\(f_{\alpha\land\beta}=\land(f_\alpha,f_\beta)\), by clone fact (i).  By
step 1 of the proof of \cref{thm:p01-w02-p01-024-nand-nor}, every
\(f\in O_n\), \(n\geq1\), is such an \(f_\varphi\).  Hence \(K=O\).

\emph{Part 3.}  If a proper clone \(K\) were included in none of the five
classes, it would contain a member outside each, and part 2 with \(C=K\)
would give \(K=\langle K\rangle=O\).  So every proper clone is included in
one of them.  Let \(X\) be one of the five and \(K\) a clone with
\(X\subsetneq K\); pick \(g\in K\setminus X\).  If \(K\neq O\), then
\(K\subseteq Y\) for one of the five classes \(Y\), and \(Y\neq X\) because
\(g\in Y\setminus X\); so \(X\subseteq Y\), contrary to part 1.  Hence
\(K=O\), and \(X\), proper by part 1, is maximal.  Conversely, a maximal
clone \(K\) is proper, hence included in some \(X\) of the five; \(X\) is a
proper clone including \(K\), so \(K=X\) by maximality.

\emph{The correspondence with formulas.}  For a set \(C\) of connectives of
positive arity, the truth functions of formulas over \(C\) in
\(p_1,\ldots,p_n\), for all \(n\geq1\), form exactly \(\langle C\rangle\).
They lie in \(\langle C\rangle\) by the computation of Step E with the
connectives of \(C\) in place of \(\neg\) and \(\land\).  Conversely, the
set \(F\) of these truth functions contains every projection
(\(f_{p_i}=e^n_i\)) and every interpreting function of \(C\)
(\(f_{c(p_1,\ldots,p_k)}=c\)), and it is closed under composition: if
\(f=f_\varphi\in O_m\) and \(g_j=f_{\psi_j}\in O_n\), then substituting
\(\psi_j\) for \(p_j\) in \(\varphi\) gives a formula whose truth function
is \(f(g_1,\ldots,g_m)\), by induction on \(\varphi\).  So \(F\) is a clone
containing \(C\), and \(\langle C\rangle\subseteq F\).  A nullary constant
symbol enters as its unary constant function: replacing \(\bottom\) by
\(k_0(p_1)\) or \(\topv\) by \(k_1(p_1)\) does not change the truth
function of a formula with at least the variable \(p_1\) available.
\end{proof}
\end{annexitem}
